\documentclass[11pt]{article}

\usepackage[margin=1in]{geometry}
\usepackage{amsmath, amssymb, amsthm}
\usepackage{booktabs}

\newtheorem{proposition}{Proposition}
\newtheorem{corollary}[proposition]{Corollary}
\newtheorem{remark}[proposition]{Remark}
\theoremstyle{definition}
\newtheorem{definition}[proposition]{Definition}

\DeclareMathOperator{\SO}{SO}
\DeclareMathOperator{\Exp}{Exp}
\DeclareMathOperator{\Log}{Log}
\DeclareMathOperator{\sinc}{sinc}
\newcommand{\so}{\mathfrak{so}}
\newcommand{\R}{\mathbb{R}}
\newcommand{\ip}[2]{\langle #1,\, #2 \rangle}
\newcommand{\skewm}[1]{[#1]_{\times}}
\newcommand{\Jl}{J_{\ell}}
\newcommand{\Jr}{J_{r}}

\title{Spatial versus chart gradients on $\SO(3)$\\
{\large why $\nabla_r f = \Jl(r)^{\!\top}\nabla_\omega f$}}
\author{}
\date{}

\begin{document}
\maketitle

\section{Setup}

Let $\skewm{\cdot}:\R^3\to\so(3)$ be the hat map,
\[
  \skewm{v} =
  \begin{pmatrix}
    0 & -v_3 & v_2\\
    v_3 & 0 & -v_1\\
    -v_2 & v_1 & 0
  \end{pmatrix},
  \qquad \skewm{v}w = v\times w,
\]
and let $\Exp:\R^3\to\SO(3)$ be the exponential map, given by Rodrigues' formula
\[
  \Exp(r) = I + \frac{\sin\theta}{\theta}\skewm{r}
          + \frac{1-\cos\theta}{\theta^{2}}\skewm{r}^{2},
  \qquad \theta := \lVert r\rVert .
\]
We call $r$ the \emph{rotation vector} (the chart coordinate) and $R=\Exp(r)$ the
group element it represents.

Let $f:\SO(3)\to\R$ be a differentiable scalar field --- in our application the
translational clearance $d$, the angular clearance $\alpha$, or the travel time
$T$. There are two inequivalent notions of ``the gradient of $f$'', and the
entire issue is that our pipeline computes one and consumes the other.

\begin{description}
\item[Chart gradient.] Set $F := f\circ\Exp:\R^3\to\R$. Then
  $\nabla_r f := \nabla F(r)\in\R^3$ is the ordinary Euclidean gradient of the
  composed map. \emph{This is what automatic differentiation returns}, because
  the network consumes $r$ (in our code, $r/2\pi$) as an input vector.

\item[Spatial gradient.] For $\omega\in\R^3$ define the left (spatial)
  directional derivative
  \[
    D_\omega f(R) := \left.\frac{d}{dt}\right|_{t=0} f\!\left(\Exp(t\omega)\,R\right).
  \]
  This is linear in $\omega$, so there is a unique $\nabla_\omega f\in\R^3$ with
  $D_\omega f = \ip{\nabla_\omega f}{\omega}$. \emph{This is what the geometric
  preprocessing computes} (see Section~\ref{sec:app}), because it perturbs the
  placement by a physical rotation, not by a coordinate increment.
\end{description}

\section{The left Jacobian}

\begin{definition}
The \emph{left Jacobian} $\Jl(r)\in\R^{3\times3}$ is defined by the first-order
expansion
\begin{equation}\label{eq:jldef}
  \Exp(r+\delta) \;=\; \Exp\!\big(\Jl(r)\,\delta + O(\lVert\delta\rVert^{2})\big)\,\Exp(r).
\end{equation}
\end{definition}

Equation~\eqref{eq:jldef} says precisely: \emph{a coordinate increment $\delta$
produces, to first order, a spatial rotation $\omega = \Jl(r)\delta$.}
It admits the closed form
\begin{equation}\label{eq:jlclosed}
  \Jl(r) \;=\; I \;+\; \frac{1-\cos\theta}{\theta^{2}}\skewm{r}
              \;+\; \frac{\theta-\sin\theta}{\theta^{3}}\skewm{r}^{2}.
\end{equation}

\section{The relation}

\begin{proposition}\label{prop:main}
For all $r$ with $\theta=\lVert r\rVert<\pi$,
\[
  \boxed{\;\nabla_r f \;=\; \Jl(r)^{\!\top}\,\nabla_\omega f\;}
\]
\end{proposition}

\begin{proof}
Fix $r$ and let $\delta\in\R^3$. Applying \eqref{eq:jldef} and then the
definition of the spatial gradient with $\omega=\Jl(r)\delta$,
\begin{align*}
  F(r+\delta)
    &= f\!\big(\Exp(r+\delta)\big) \\
    &= f\!\left(\Exp\big(\Jl(r)\delta + O(\lVert\delta\rVert^{2})\big)\,\Exp(r)\right)\\
    &= f\!\big(\Exp(r)\big) \;+\; \ip{\nabla_\omega f}{\Jl(r)\delta} \;+\; O(\lVert\delta\rVert^{2})\\
    &= F(r) \;+\; \ip{\Jl(r)^{\!\top}\nabla_\omega f}{\delta} \;+\; O(\lVert\delta\rVert^{2}).
\end{align*}
Comparing with $F(r+\delta) = F(r) + \ip{\nabla F(r)}{\delta} + O(\lVert\delta\rVert^2)$
and using uniqueness of the gradient gives $\nabla F(r) = \Jl(r)^{\!\top}\nabla_\omega f$.
\end{proof}

\begin{remark}[Why the transpose]
A gradient is a \emph{covector}. Tangent vectors push forward by $\Jl$
(Eq.~\eqref{eq:jldef}); covectors therefore pull back by $\Jl^{\!\top}$. Confusing
the two is the practical trap: see Corollary~\ref{cor:vecvscovec}.
\end{remark}

\section{Structure of $\Jl$}

Write $\hat r = r/\theta$. Using $\skewm{r}=\theta\skewm{\hat r}$ and
$\skewm{r}^{2}=\theta^{2}\big(\hat r\hat r^{\top}-I\big)$ in \eqref{eq:jlclosed}
and collecting terms,
\begin{equation}\label{eq:jlspectral}
  \Jl(r) \;=\; \frac{\sin\theta}{\theta}\,I
         \;+\; \left(1-\frac{\sin\theta}{\theta}\right)\hat r\hat r^{\top}
         \;+\; \frac{1-\cos\theta}{\theta}\,\skewm{\hat r}.
\end{equation}

\begin{proposition}\label{prop:spectral}
Decompose $\R^{3}=\operatorname{span}(\hat r)\oplus \hat r^{\perp}$. Then
\begin{enumerate}
\item $\Jl(r)\,\hat r = \hat r$: the axial direction is fixed exactly, with unit gain.
\item On $\hat r^{\perp}$, $\Jl(r)$ acts as a rotation by $\theta/2$ within that
  plane, scaled by $\sinc(\theta/2) = \dfrac{\sin(\theta/2)}{\theta/2}$.
\end{enumerate}
\end{proposition}

\begin{proof}
(1) is immediate from \eqref{eq:jlspectral}, since $\hat r\hat r^\top \hat r=\hat r$
and $\skewm{\hat r}\hat r=0$, giving
$\Jl\hat r = \tfrac{\sin\theta}{\theta}\hat r + (1-\tfrac{\sin\theta}{\theta})\hat r = \hat r$.

For (2), on $\hat r^{\perp}$ the term $\hat r\hat r^{\top}$ vanishes and
$\skewm{\hat r}$ acts as a rotation by $\pi/2$ inside the plane. Hence
$\Jl|_{\perp} = a\,I_{\perp} + b\,\mathrm{Rot}_{\pi/2}$ with
$a=\sin\theta/\theta$, $b=(1-\cos\theta)/\theta$, which is the rotation by
$\arctan(b/a)$ scaled by $\sqrt{a^{2}+b^{2}}$. Now
\[
  a^{2}+b^{2}
  = \frac{\sin^{2}\theta + 1 - 2\cos\theta + \cos^{2}\theta}{\theta^{2}}
  = \frac{2-2\cos\theta}{\theta^{2}}
  = \frac{4\sin^{2}(\theta/2)}{\theta^{2}},
\]
so $\sqrt{a^2+b^2} = \dfrac{2\sin(\theta/2)}{\theta} = \sinc(\theta/2)$, and
\[
  \arctan\frac{b}{a}
  = \arctan\frac{1-\cos\theta}{\sin\theta}
  = \arctan\frac{2\sin^{2}(\theta/2)}{2\sin(\theta/2)\cos(\theta/2)}
  = \arctan\tan\frac{\theta}{2}
  = \frac{\theta}{2}. \qedhere
\]
\end{proof}

\begin{corollary}
$\Jl(r)^{\!\top} = \Jl(-r) =: \Jr(r)$, the right Jacobian. It fixes $\hat r$ and
acts on $\hat r^{\perp}$ as $\sinc(\theta/2)$ times a rotation by $-\theta/2$.
\end{corollary}

\begin{corollary}[Combined probe attenuation]
Take a coordinate step $\delta = s\,u$ with $u\perp\hat r$ a unit vector. The
induced spatial rotation has magnitude $s\,\sinc(\theta/2)$, and its axis is
rotated by $\theta/2$ away from $u$. The component of the rotation actually
about $u$ therefore carries the factor
\[
  \sinc(\theta/2)\cos(\theta/2)
  = \frac{2\sin(\theta/2)\cos(\theta/2)}{\theta}
  = \frac{\sin\theta}{\theta} = \sinc(\theta),
\]
which vanishes at $\theta=\pi$.
\end{corollary}

\begin{corollary}[Vectors versus covectors]\label{cor:vecvscovec}
Let $n$ be a direction expressed in the spatial frame.
\begin{itemize}
\item To \emph{step} so that the placement rotates about $n$ (a tangent vector),
  the chart increment is $\delta \propto \Jl(r)^{-1} n$.
\item To \emph{compare} $n$ against a chart-space gradient such as
  $\nabla_r T$ (a covector), the chart-space direction is
  $\propto \Jl(r)^{\!\top} n$.
\end{itemize}
These differ. Since $\Jl^{-1}$ and $\Jl^{\!\top}$ both fix $\hat r$ and both
rotate $\hat r^{\perp}$ by $-\theta/2$ but with reciprocal scalings
$\sinc(\theta/2)^{-1}$ and $\sinc(\theta/2)$, they act identically only on
vectors that are purely axial or purely perpendicular.
\end{corollary}

\section{Application to the preprocessed normals}\label{sec:app}

Both geometric quantities produced by the preprocessing are \emph{spatial}
gradients, so both require Proposition~\ref{prop:main} before being compared
with a network gradient.

\paragraph{Translational clearance.}
Let $s^{\star}$ be the closest point of the placed body to the obstacle set,
$c$ the placement centre, $\rho := s^{\star}-c$, and $n_{ws}$ the unit contact
normal. An infinitesimal spatial rotation $\omega$ displaces $s^{\star}$ by
$\omega\times\rho$, so the clearance $d$ changes by
\[
  \delta d = \ip{n_{ws}}{\omega\times\rho} = \ip{\omega}{\rho\times n_{ws}}
  \quad\Longrightarrow\quad
  \nabla_\omega d = \rho\times n_{ws} \;=:\; \mathrm{moment}.
\]
Hence the chart gradient, in the $2\pi$-normalised coordinate
$\tilde r = r/2\pi$, is
\[
  \nabla_{\tilde r}\, d \;=\; 2\pi\,\Jl(r)^{\!\top}\big(\rho\times n_{ws}\big),
\]
whereas the code stores $2\pi\,(\rho\times n_{ws})$, i.e.\ it omits
$\Jl(r)^{\!\top}$.

\paragraph{Angular clearance.}
Let $\alpha=\arccos\ip{u_s}{u_e}$ with $u_s,u_e$ unit spatial directions from
the centre. Then $\delta\ip{u_s}{u_e} = \ip{\omega\times u_s}{u_e}
= \ip{\omega}{u_s\times u_e}$, and since
$\delta\alpha = -\delta\cos\alpha/\sin\alpha$,
\[
  \nabla_\omega \alpha \;=\; -\,\frac{u_s\times u_e}{\sin\alpha}.
\]
So the stored axis $\texttt{rot\_n}=\widehat{u_s\times u_e}$ is the normalised
spatial direction of steepest \emph{decrease} of $\alpha$. Its chart-space
counterpart --- the object that may legitimately be aligned against
$\nabla_{\tilde r}T$ --- is
\[
  \widehat{\Jl(r)^{\!\top}\big(u_s\times u_e\big)}.
\]

\begin{remark}[The $2\pi$ is harmless; the Jacobian is not]
Rescaling $r\mapsto r/2\pi$ multiplies the gradient by the scalar $2\pi$ and
therefore leaves \emph{directions} unchanged. Any purely directional loss (e.g.\
normal alignment) is insensitive to it. By contrast $\Jl^{\!\top}$ scales the
component perpendicular to $\hat r$ by $\sinc(\theta/2)$ while fixing the axial
component, \emph{and} rotates the perpendicular part by $-\theta/2$; both change
the direction.
\end{remark}

\paragraph{Magnitude of the error.}
Table~\ref{tab:err} reports the angle between $n$ and
$\widehat{\Jl(r)^{\!\top}n}$ for $n$ drawn uniformly on the sphere and
independent of $\hat r$, which matches how the contact axis $u_s\times u_e$
relates to the base rotation in our sampler.

\begin{table}[h]
\centering
\begin{tabular}{lrr}
\toprule
$\theta=\lVert r\rVert$ & mean dev. & max dev. \\
\midrule
$0.25$   & $5.6^{\circ}$  & $7.2^{\circ}$ \\
$0.50$   & $11.3^{\circ}$ & $14.3^{\circ}$ \\
$1.00$   & $22.5^{\circ}$ & $28.6^{\circ}$ \\
$\pi/2$  & $34.9^{\circ}$ & $45.0^{\circ}$ \\
$2.00$   & $43.8^{\circ}$ & $57.3^{\circ}$ \\
$2.50$   & $53.5^{\circ}$ & $71.6^{\circ}$ \\
$3.00$   & $63.0^{\circ}$ & $85.9^{\circ}$ \\
\midrule
$\theta\sim\mathcal{U}[0,\pi]$ & \multicolumn{2}{c}{mean $33.7^{\circ}$,\; median $30.6^{\circ}$,\; p95 $75.0^{\circ}$} \\
\bottomrule
\end{tabular}
\caption{Misalignment incurred by treating the spatial axis as a chart-space
covector. The base rotation angle is drawn as $\theta\sim\mathcal{U}[0,\pi]$ by
the sampler, giving the aggregate figures in the last row.}
\label{tab:err}
\end{table}

\begin{remark}[Translation]
For the translational block the configuration space is the vector space $\R^{3}$:
the chart is the group, $\Jl\equiv I$, and vectors and covectors transform
identically. The stored $n_{ws}$ therefore needs no correction, which is why the
translational terms behave while the rotational ones do not.
\end{remark}

\end{document}
